% GENERATED FILE. Edit canonical lessons, then run npm run generate:books.
% canonical-source-hash: fnv1a64:8b0f9632edd6dbfc
% canonical-lessons: MW-C187-kotho, MW-C187-alo, MW-C187-bari, MW-C187-khuno, MW-C187-maliyo

\chapter{Rooms and Corners}
\label{ch:mw-rooms-and-corners}
\hlchaptermodality{187}

\begin{chapteropening}
\textbf{By the end of this chapter:} \emph{I can name the rooms and corners of a house.}

Complete the last lesson of chapter 187: I can name the rooms and corners of a house.
\end{chapteropening}

\section[\texorpdfstring{\mw{कोठो} (koṭho)}{कोठो (koṭho)}]{\mw{कोठो} (koṭho) — a storeroom}

\label{lesson:MW-C187-kotho}



\begin{quote}
Before the new one: say the Marwadi for a plank seat, then the Marwadi for a bedstead.
\end{quote}

\subsection*{You'll want to know: \mw{कोठो}}
\textbf{\mw{कोठो}} — \emph{koṭho} — \textquotedblleft{}a storeroom\textquotedblright{}.

\begin{quote}
\mw{म्हारो} \mw{कोठो} \mw{छोटो} \mw{है।}
\end{quote}

\begin{quote}
\emph{mhāro koṭho chhoṭo hai.}
\end{quote}

\begin{quote}
My storeroom is small.
\end{quote}

A masculine noun.

\begin{cousinweb}
\textbf{\mw{कोठो}} (\emph{koṭho}) is inherited from Sanskrit \emph{koṣṭha}, \textquotedblleft{}a granary, a storeroom\textquotedblright{}; CDIAL lists Marwari \emph{koṭho}, \textquotedblleft{}a room\textquotedblright{} (CDIAL 3546).
\end{cousinweb}

\begin{grammarlens}[title={What you've built}]
One more word for this chapter.
\end{grammarlens}

\subsection*{Your turn}
\begin{itemize}
\raggedright
  \item \emph{Say it:} \emph{koṭho}
  \item \emph{Say it:} \emph{koṭho}, once more
  \item \emph{Say it:} \emph{mhāro koṭho chhoṭo hai.} — My storeroom is small.
  \item \emph{Recall:} say the Marwadi for a plank seat, then the Marwadi for a bedstead, then say \emph{koṭho} again
\end{itemize}

\begin{tcolorbox}[breakable,colback=teal!4,colframe=teal!35!black,title={Before you move on}]
What does \mw{कोठो} mean? (\textquotedblleft{}A storeroom\textquotedblright{}.) Say it once more.
\end{tcolorbox}
\section[\texorpdfstring{\mw{आळो} (āḷo)}{आळो (āḷo)}]{\mw{आळो} (āḷo) — a wall niche}

\label{lesson:MW-C187-alo}



\begin{quote}
Before the new one: say the Marwadi for a bedstead, then the Marwadi for a storeroom.
\end{quote}

\subsection*{You'll want to know: \mw{आळो}}
\textbf{\mw{आळो}} — \emph{āḷo} — \textquotedblleft{}a wall niche\textquotedblright{}.

\begin{quote}
\mw{भींत} \mw{मांय} \mw{एक} \mw{आळो} \mw{है।}
\end{quote}

\begin{quote}
\emph{bhī̃t mā̃y ek āḷo hai.}
\end{quote}

\begin{quote}
There's a niche in the wall.
\end{quote}

A masculine noun.

\begin{cousinweb}
\textbf{\mw{आळो}} (\emph{āḷo}) is inherited from Sanskrit \emph{ālaya}, \textquotedblleft{}a house, a receptacle\textquotedblright{}; CDIAL lists Hindi \emph{ālā}, \textquotedblleft{}a niche in a wall\textquotedblright{} (CDIAL 1366). CDIAL gives no Marwari form in this entry.
\end{cousinweb}

\begin{grammarlens}[title={What you've built}]
One more word for this chapter.
\end{grammarlens}

\subsection*{Your turn}
\begin{itemize}
\raggedright
  \item \emph{Say it:} \emph{āḷo}
  \item \emph{Say it:} \emph{āḷo}, once more
  \item \emph{Say it:} \emph{bhī̃t mā̃y ek āḷo hai.} — There's a niche in the wall.
  \item \emph{Recall:} say the Marwadi for a bedstead, then the Marwadi for a storeroom, then say \emph{āḷo} again
\end{itemize}

\begin{tcolorbox}[breakable,colback=teal!4,colframe=teal!35!black,title={Before you move on}]
What does \mw{आळो} mean? (\textquotedblleft{}A wall niche\textquotedblright{}.) Say it once more.
\end{tcolorbox}
\section[\texorpdfstring{\mw{बारी} (bārī)}{बारी (bārī)}]{\mw{बारी} (bārī) — a small window}

\label{lesson:MW-C187-bari}



\begin{quote}
Before the new one: say the Marwadi for a storeroom, then the Marwadi for a wall niche.
\end{quote}

\subsection*{You'll want to know: \mw{बारी}}
\textbf{\mw{बारी}} — \emph{bārī} — \textquotedblleft{}a small window\textquotedblright{}.

\begin{quote}
\mw{रसोई} \mw{मांय} \mw{एक} \mw{छोटी} \mw{बारी} \mw{है।}
\end{quote}

\begin{quote}
\emph{rasoī mā̃y ek chhoṭī bārī hai.}
\end{quote}

\begin{quote}
There's a small window in the kitchen.
\end{quote}

A feminine noun.

\begin{cousinweb}
\textbf{\mw{बारी}} (\emph{bārī}) is inherited from Sanskrit \emph{dvāra}, \textquotedblleft{}a door\textquotedblright{}; CDIAL lists Old Marwari \emph{bāra}, \textquotedblleft{}a door\textquotedblright{}, and Gujarati \emph{bārī}, \textquotedblleft{}a window\textquotedblright{} (CDIAL 6663).
\end{cousinweb}

\begin{grammarlens}[title={What you've built}]
One more word for this chapter.
\end{grammarlens}

\subsection*{Your turn}
\begin{itemize}
\raggedright
  \item \emph{Say it:} \emph{bārī}
  \item \emph{Say it:} \emph{bārī}, once more
  \item \emph{Say it:} \emph{rasoī mā̃y ek chhoṭī bārī hai.} — There's a small window in the kitchen.
  \item \emph{Recall:} say the Marwadi for a storeroom, then the Marwadi for a wall niche, then say \emph{bārī} again
\end{itemize}

\begin{tcolorbox}[breakable,colback=teal!4,colframe=teal!35!black,title={Before you move on}]
What does \mw{बारी} mean? (\textquotedblleft{}A small window\textquotedblright{}.) Say it once more.
\end{tcolorbox}
\section[\texorpdfstring{\mw{खूणो} (khūṇo)}{खूणो (khūṇo)}]{\mw{खूणो} (khūṇo) — a corner}

\label{lesson:MW-C187-khuno}



\begin{quote}
Before the new one: say the Marwadi for a wall niche, then the Marwadi for a small window.
\end{quote}

\subsection*{You'll want to know: \mw{खूणो}}
\textbf{\mw{खूणो}} — \emph{khūṇo} — \textquotedblleft{}a corner\textquotedblright{}.

\begin{quote}
\mw{खूणो} \mw{खाली} \mw{है।}
\end{quote}

\begin{quote}
\emph{khūṇo khālī hai.}
\end{quote}

\begin{quote}
The corner is empty.
\end{quote}

A masculine noun.

\begin{cousinweb}
\textbf{\mw{खूणो}} (\emph{khūṇo}) comes from a reconstructed \emph{*khuṇṭa}, \textquotedblleft{}a corner\textquotedblright{}; CDIAL lists Gujarati \emph{khūṇo} and Hindi \emph{khū̃ṭ} (CDIAL 3898). CDIAL gives no Marwari form in this entry.
\end{cousinweb}

\begin{grammarlens}[title={What you've built}]
One more word for this chapter.
\end{grammarlens}

\subsection*{Your turn}
\begin{itemize}
\raggedright
  \item \emph{Say it:} \emph{khūṇo}
  \item \emph{Say it:} \emph{khūṇo}, once more
  \item \emph{Say it:} \emph{khūṇo khālī hai.} — The corner is empty.
  \item \emph{Recall:} say the Marwadi for a wall niche, then the Marwadi for a small window, then say \emph{khūṇo} again
\end{itemize}

\begin{tcolorbox}[breakable,colback=teal!4,colframe=teal!35!black,title={Before you move on}]
What does \mw{खूणो} mean? (\textquotedblleft{}A corner\textquotedblright{}.) Say it once more.
\end{tcolorbox}
\section[\texorpdfstring{\mw{माळियो} (māḷiyo)}{माळियो (māḷiyo)}]{\mw{माळियो} (māḷiyo) — a loft}

\label{lesson:MW-C187-maliyo}



\begin{quote}
Before the new one: say the Marwadi for a small window, then the Marwadi for a corner.
\end{quote}

\subsection*{You'll want to know: \mw{माळियो}}
\textbf{\mw{माळियो}} — \emph{māḷiyo} — \textquotedblleft{}a loft\textquotedblright{}.

\begin{quote}
\mw{माळियो} \mw{खाली} \mw{कोनी।}
\end{quote}

\begin{quote}
\emph{māḷiyo khālī konī.}
\end{quote}

\begin{quote}
The loft isn't empty.
\end{quote}

A masculine noun.

\begin{cousinweb}
\textbf{\mw{माळियो}} (\emph{māḷiyo}) comes from a reconstructed \emph{*māḍa}, \textquotedblleft{}the upper storey of a house\textquotedblright{}; CDIAL lists Gujarati \emph{māḷiyũ}, \textquotedblleft{}a loft\textquotedblright{} (CDIAL 9996). CDIAL gives no Marwari form in this entry.
\end{cousinweb}

\begin{grammarlens}[title={What you've built}]
One more word for this chapter.
\end{grammarlens}

\subsection*{Your turn}
\begin{itemize}
\raggedright
  \item \emph{Say it:} \emph{māḷiyo}
  \item \emph{Say it:} \emph{māḷiyo}, once more
  \item \emph{Say it:} \emph{māḷiyo khālī konī.} — The loft isn't empty.
  \item \emph{Recall:} say the Marwadi for a small window, then the Marwadi for a corner, then say \emph{māḷiyo} again
\end{itemize}

\begin{tcolorbox}[breakable,colback=teal!4,colframe=teal!35!black,title={Before you move on}]
What does \mw{माळियो} mean? (\textquotedblleft{}A loft\textquotedblright{}.) Say it once more.
\end{tcolorbox}
